\section{Introduction}
\label{sec:intro}

BoardGameGeek's community rating aggregates the satisfaction judgments of hundreds of
thousands of players, making it the closest available proxy for long-term player
experience in the board game domain. Can the design profile of a game --- its position
in TIGER space --- predict that satisfaction score? If yes, TIGER profiles are not
merely descriptive but predictive, and the framework gains practical utility for
publishers and designers evaluating games before release.

\subsection{The BGG Rating as a Prediction Target}

BGG ratings are noisy, influenced by recency bias, publisher marketing, and
community demographics, but they represent the largest available dataset of player
experience judgments for board games~\cite{booth2021game}. Prior work has shown that
BGG weight (perceived complexity) correlates with rating within specific genres but
not across the full catalog~\cite{woods2012eurogames}. This paper tests whether
TIGER's five-dimensional profile provides a richer predictor than weight alone.

\subsection{Hypothesis}

We hypothesize that TIGER profiles predict BGG ratings at $r \geq 0.60$ correlation
on a held-out test set, representing $\geq 15\%$ RMSE improvement over the weight-only
baseline. This threshold is set to exceed the weight-only baseline by a margin that
would be practically meaningful for pre-release quality forecasting --- not merely
statistically significant.
